%% kapa: exam=1 round=28 year=2017 period=2 form=A subject=appraisal_law subject_ko=감정평가관계법규 seq=3 number=3
%% source: Q-Net 감정평가사 기출문제 (cst00309) artlSeq=5199603 file=extracted/1008/5199603/2017년 제28회 감정평가사 1차시험 문제지/2교시 A형.pdf
%% conversion: kordoc 4.24.51 -> tools/gen1.py (auto); review=proofread-sonnet; flags=none
\begin{question}{3}
국토의 계획 및 이용에 관한 법령상 용도지역의 세분에 관한 내용으로 옳게 연결한 것은?
\begin{choices}
  \item 제2종전용주거지역 - 단독주택 중심의 양호한 주거환경을 보호하기 위하여 필요한 지역
  \item 보전녹지지역 - 주로 농업적 생산을 위하여 개발을 유보할 필요가 있는 지역
  \item 제3종일반주거지역 - 중고층주택을 중심으로 편리한 주거환경을 조성하기 위하여 필요한 지역
  \item 일반상업지역 - 도심ㆍ부도심의 상업기능 및 업무기능의 확충을 위하여 필요한 지역
  \item 전용공업지역 - 경공업 그 밖의 공업을 수용하되, 주거기능ㆍ상업기능 및 업무기능의 보완이 필요한 지역
\end{choices}
\end{question}
\answer{3}
